\needspace{6\baselineskip}
\kactlsection{template}{01a581}
\begin{lstlisting}[language=C++]
#include <bits/stdc++.h>
#include <cassert>
using namespace std;

using ll = long long;
using vi = vector<int>;
using vvi = vector<vi>;
using pii = pair<int, int>;
using vpi = vector<pii>;
using vll = vector<ll>;
using i128 = __int128_t;
using ld = long double;

#define pb push_back
#define all(x) (x).begin(), (x).end()
#define sz(c) int((c).size())
#define LSOne(S) ((S) & -(S))

#define rep2(i, n) for (int i = 0; i < (n); ++i)
#define rep3(i, a, b) for (int i = (a); i < (b); ++i)
#define GET_REP(_1, _2, _3, NAME, ...) NAME
#define rep(...) GET_REP(__VA_ARGS__, rep3, rep2)(__VA_ARGS__)
#define FOR(i, a, b) for (int i = (a); i <= (b); ++i)
#define ROF(i, a, b) for (int i = (a); i >= (b); --i)

const int INF = 1e9 + 7;
const ll LINF = 1e18 + 7;
const int MOD = 1e9 + 7;
const double EPS = 1e-9;

// up, down, left, right
const int dx[] = {-1, 1, 0, 0};
const int dy[] = {0, 0, -1, 1};

ll binpow(ll a, ll b, ll mod) {
  ll res = 1;
  a %= mod;
  while (b) {
    if (b & 1)
      res = res * a % mod;
    a = a * a % mod;
    b >>= 1;
  }
  return res;
}

void solve() {}

int main() {
  ios_base::sync_with_stdio(false);
  cin.tie(nullptr);
  solve();
  return 0;
}
\end{lstlisting}
